\documentclass[lualatex,ja=standard,a4paper,11pt]{jlreq}

\usepackage{amsmath,amssymb,amsfonts}
\usepackage{luatexja}
\usepackage{tikz}
\usepackage{geometry}
\geometry{top=20mm,bottom=20mm,left=20mm,right=20mm}

\title{ {{title}} }
\author{高校数学教材AI v20}
\date{\today}

\begin{document}

\maketitle

\begin{center}
\textbf{【注意事項】} \\
1. 問題は全部で {{ problems|length }} 問あります。 \\
2. 各問題の配点は問題文の末尾に記載されています。
\end{center}

\section*{問題}

{% for item in problems %}
\subsection*{第 {{ loop.index }} 問 （{{ item.points }}点）}
{{ item.problem }}

{% if item.figure_code %}
\begin{center}
{{ item.figure_code }}
\end{center}
{% endif %}

\vspace{30mm} % 解答欄のスペース
{% endfor %}

\newpage
\section*{解答・解説}

{% for item in problems %}
\subsection*{第 {{ loop.index }} 問 の解答}
\textbf{【答え】} {{ item.answer_key }}

\vspace{5mm}
\textbf{【解説】}
{{ item.solution }}

\vspace{10mm}
{% endfor %}

\end{document}
